\documentclass[12pt,a4paper]{article}

% ===== PACKAGES =====
\usepackage[utf8]{inputenc}
\usepackage[vietnamese]{babel}
\usepackage{amsmath,amssymb}
\usepackage{graphicx}
\usepackage{booktabs}
\usepackage{multirow}
\usepackage{tabularx}
\usepackage{xcolor}
\usepackage{colortbl}
\usepackage{hyperref}
\usepackage{enumitem}
\usepackage{geometry}
\usepackage{fancyhdr}
\usepackage{tikz}
\usetikzlibrary{shapes.geometric, arrows.meta, positioning, calc, decorations.pathreplacing}
\usepackage{tcolorbox}
\tcbuselibrary{skins,breakable}
\usepackage{float}
\usepackage{titlesec}
\usepackage{pgfgantt}

% ===== PAGE SETUP =====
\geometry{left=2.5cm, right=2.5cm, top=2.5cm, bottom=2.5cm}
\setlength{\parindent}{1.2em}
\setlength{\parskip}{0.4em}

% ===== COLORS =====
\definecolor{primary}{HTML}{1B5E20}
\definecolor{secondary}{HTML}{2E7D32}
\definecolor{accent}{HTML}{00695C}
\definecolor{highlight}{HTML}{E65100}
\definecolor{lightbg}{HTML}{E8F5E9}
\definecolor{lightaccent}{HTML}{E0F2F1}
\definecolor{tableheader}{HTML}{1B5E20}
\definecolor{tablerow1}{HTML}{F5F5F5}
\definecolor{tablerow2}{HTML}{FFFFFF}
\definecolor{phase1}{HTML}{1565C0}
\definecolor{phase2}{HTML}{6A1B9A}
\definecolor{phase3}{HTML}{E65100}
\definecolor{phase4}{HTML}{2E7D32}
\definecolor{done}{HTML}{4CAF50}
\definecolor{inprog}{HTML}{FF9800}
\definecolor{todo}{HTML}{9E9E9E}

% ===== HEADER/FOOTER =====
\pagestyle{fancy}
\fancyhf{}
\fancyhead[L]{\small\color{primary}\textit{Kế Hoạch Nghiên Cứu — GEDA}}
\fancyhead[R]{\small\color{primary}\textit{Natural Disaster Classification}}
\fancyfoot[C]{\thepage}
\renewcommand{\headrulewidth}{0.5pt}
\renewcommand{\headrule}{\hbox to\headwidth{\color{primary}\leaders\hrule height \headrulewidth\hfill}}

% ===== SECTION STYLING =====
\titleformat{\section}{\Large\bfseries\color{primary}}{\thesection}{1em}{}[\color{primary}\titlerule]
\titleformat{\subsection}{\large\bfseries\color{secondary}}{\thesubsection}{1em}{}
\titleformat{\subsubsection}{\normalsize\bfseries\color{accent}}{\thesubsubsection}{1em}{}

% ===== TCOLORBOX STYLES =====
\newtcolorbox{goalbox}[1][]{
  enhanced, breakable,
  colback=lightbg, colframe=primary,
  fonttitle=\bfseries\color{white},
  title=#1, coltitle=white, colbacktitle=primary,
  boxrule=0.8pt, arc=3pt,
  left=8pt, right=8pt, top=6pt, bottom=6pt
}

\newtcolorbox{taskbox}[1][]{
  enhanced, breakable,
  colback=lightaccent, colframe=accent,
  fonttitle=\bfseries\color{white},
  title=#1, coltitle=white, colbacktitle=accent,
  boxrule=0.8pt, arc=3pt,
  left=8pt, right=8pt, top=6pt, bottom=6pt
}

\newtcolorbox{riskbox}[1][]{
  enhanced, breakable,
  colback=orange!5, colframe=highlight,
  fonttitle=\bfseries\color{white},
  title=#1, coltitle=white, colbacktitle=highlight,
  boxrule=0.8pt, arc=3pt,
  left=8pt, right=8pt, top=6pt, bottom=6pt
}

% ===== HYPERREF =====
\hypersetup{
  colorlinks=true, linkcolor=primary, urlcolor=accent, citecolor=secondary,
  pdftitle={GEDA Research Project Plan},
  pdfauthor={CrisisSpot Research Team}
}

\begin{document}

% ========== TITLE PAGE ==========
\begin{titlepage}
\begin{center}
\vspace*{1cm}

\begin{tikzpicture}
  \fill[primary] (0,0) rectangle (\textwidth, 0.4);
\end{tikzpicture}

\vspace{1.5cm}
{\Huge\bfseries\color{primary} Kế Hoạch Nghiên Cứu}

\vspace{0.6cm}
{\Large\color{secondary} Graph-Enhanced Differential Attention\\cho Phân Loại Thảm Hoạ Thiên Nhiên\\Đa Phương Thức trên Mạng Xã Hội}

\vspace{1cm}
\begin{tikzpicture}
  \fill[accent] (0,0) rectangle (8cm, 0.15);
\end{tikzpicture}

\vspace{1cm}
{\large\color{primary}\textbf{Tên dự án:} GEDA — Graph-Enhanced Differential Attention}\\[0.4cm]
{\normalsize
\textbf{Lĩnh vực:} Multimodal Machine Learning $\times$ Natural Disaster Response\\[0.2cm]
\textbf{Dataset:} CrisisMMD v2.0 (7 natural disasters, 20K tweets)\\[0.2cm]
\textbf{Stack:} Python $\cdot$ PyTorch $\cdot$ CLIP $\cdot$ GraphSAGE $\cdot$ FAISS $\cdot$ Google Colab\\[0.2cm]
\textbf{Dựa trên:} IEEE Access 2025 (GNN) + CVPRw MMFM 2025 (DiffAttn)
}

\vspace{1.2cm}
{\large Dự án: \textbf{CrisisSummarization}}\\[0.3cm]
{\normalsize Ngày lập kế hoạch: 04/03/2026}

\vfill
\begin{tikzpicture}
  \fill[primary] (0,0) rectangle (\textwidth, 0.4);
\end{tikzpicture}
\end{center}
\end{titlepage}

% ========== TOC ==========
\tableofcontents
\newpage

% ================================================================
% SECTION 1: TỔNG QUAN DỰ ÁN
% ================================================================
\section{Tổng Quan Dự Án}

\subsection{Tầm Nhìn}

\begin{goalbox}[Tầm nhìn dự án]
Xây dựng hệ thống \textbf{phân loại thảm hoạ thiên nhiên đa phương thức} trên mạng xã hội, kết hợp sức mạnh của \textbf{Graph Neural Networks} (học từ cấu trúc dữ liệu) và \textbf{Differential Attention} (lọc nhiễu thích ứng), đạt SOTA trên CrisisMMD v2.0 trong điều kiện few-shot, hướng tới ứng dụng thực tế hỗ trợ ứng phó thảm hoạ.
\end{goalbox}

\subsection{Phạm Vi Nghiên Cứu}

\begin{table}[H]
\centering
\caption{Phạm vi nghiên cứu}
\rowcolors{2}{tablerow1}{tablerow2}
\begin{tabularx}{\textwidth}{>{\bfseries}l X}
\toprule
\rowcolor{tableheader}
\textcolor{white}{\textbf{Hạng mục}} & \textcolor{white}{\textbf{Chi tiết}} \\
\midrule
Loại thảm hoạ & Thiên nhiên: bão, động đất, lũ lụt, cháy rừng \\
Dataset & CrisisMMD v2.0 (7 events, 20K tweets) \\
Tasks & Task 1: Informativeness $\cdot$ Task 2: Humanitarian $\cdot$ Task 3: Damage Severity \\
Few-shot settings & 50, 100, 250, 500 labeled samples \\
Modalities & Text (tweets) + Image (photos) \\
Ngoài phạm vi & Biểu tình chính trị (7Sept), dịch bệnh, tấn công mạng \\
\bottomrule
\end{tabularx}
\end{table}

\subsection{Mục Tiêu Cụ Thể}

\begin{enumerate}[leftmargin=*]
  \item \textbf{MT1:} Xây dựng kiến trúc GEDA kết hợp GNN + Attention, đạt F1 cao hơn cả 2 baselines trên ít nhất 2/3 tasks
  \item \textbf{MT2:} Giải quyết scalability ($O(n^2) \rightarrow O(n \log n)$) bằng FAISS
  \item \textbf{MT3:} Thiết kế Multi-task Learning pipeline train cả 3 tasks đồng thời
  \item \textbf{MT4:} Đánh giá cross-disaster-type generalization (leave-one-type-out)
  \item \textbf{MT5:} Đảm bảo statistical rigor: 5 seeds, CI, significance tests cho mọi kết quả
\end{enumerate}

% ================================================================
% SECTION 2: KIẾN TRÚC ĐỀ XUẤT
% ================================================================
\section{Kiến Trúc GEDA}

\subsection{Sơ Đồ Tổng Thể}

\begin{figure}[H]
\centering
\begin{tikzpicture}[
  node distance=0.85cm,
  block/.style={rectangle, draw=#1, fill=#1!8, text width=7cm, minimum height=0.85cm, align=center, rounded corners=2pt, font=\footnotesize\bfseries},
  phase/.style={font=\tiny\color{gray!70}, align=left},
  arrow/.style={-{Stealth[length=2.5mm]}, thick, color=primary!70}
]
  % Input
  \node[block=primary] (input) {Social Media Post: Image $I$ + Text $T$};
  
  % Stage 1: Knowledge Fusion
  \node[block=phase2, below=of input] (kf) {Stage 1: LLaVA Knowledge Fusion};
  \node[phase, left=0.3cm of kf] {từ Paper 2};
  
  % Stage 2: Feature Extraction
  \node[block=phase1, below=of kf] (clip) {Stage 2: CLIP Feature Extraction (frozen)};
  \node[phase, left=0.3cm of clip] {chung};

  % Stage 3: Dimensionality Reduction
  \node[block=phase1, below=of clip] (pca) {Stage 3: PCA Reduction $\rightarrow$ 256d per modality};
  \node[phase, left=0.3cm of pca] {từ Paper 1};
  
  % Stage 4: Graph Construction
  \node[block=phase3, below=of pca] {Stage 4: FAISS k-NN Graph Construction ($k$=16)};
  \node[phase, left=0.3cm of pca.south west, yshift=-0.85cm] (s4l) {\textcolor{highlight}{\textbf{MỚI}}};
  
  % Stage 5: Graph Learning
  \node[block=phase1, below=2cm of pca] (gnn) {Stage 5: GraphSAGE (2 layers $\times$ 512d)};
  \node[phase, left=0.3cm of gnn] {từ Paper 1};
  
  % Stage 6: Attention Stack
  \node[block=phase2, below=of gnn] (attn) {Stage 6: Self-Attn $\rightarrow$ GCA $\rightarrow$ Adaptive DiffAttn};
  \node[phase, left=0.3cm of attn] {từ Paper 2};
  
  % Stage 7: Multi-task Head
  \node[block=phase3, below=of attn] (mtl) {Stage 7: Multi-task Classification (Task 1 + 2 + 3)};
  \node[phase, left=0.3cm of mtl] {\textcolor{highlight}{\textbf{MỚI}}};
  
  % Arrows
  \draw[arrow] (input) -- (kf);
  \draw[arrow] (kf) -- (clip);
  \draw[arrow] (clip) -- (pca);
  \draw[arrow] (pca) -- ++(0,-0.85);
  \draw[arrow] ([yshift=0.42cm]gnn.north) -- (gnn);
  \draw[arrow] (gnn) -- (attn);
  \draw[arrow] (attn) -- (mtl);
  
  % Brace for "Novel"
  \draw[decorate,decoration={brace,amplitude=8pt,mirror},thick,primary] 
    ([xshift=4.5cm]gnn.north east) -- ([xshift=4.5cm]mtl.south east)
    node[midway,right=10pt,font=\small\bfseries\color{primary}] {GEDA Core};
\end{tikzpicture}
\caption{Kiến trúc GEDA đề xuất — 7 stages}
\label{fig:geda-full}
\end{figure}

\subsection{Loss Function}

\begin{equation}
\boxed{
\mathcal{L}_{\text{GEDA}} = \underbrace{\sum_{t=1}^{3} \lambda_t \cdot \text{FocalLoss}_t}_{\text{Multi-task}} + \underbrace{\alpha \cdot \mathcal{L}_{\text{contrastive}}}_{\text{Shared repr.}} + \underbrace{\beta \cdot \mathcal{L}_{\text{graph}}}_{\text{Graph consistency}}
}
\end{equation}

\begin{table}[H]
\centering
\caption{Tham số Loss Function}
\rowcolors{2}{tablerow1}{tablerow2}
\begin{tabular}{llll}
\toprule
\rowcolor{tableheader}
\textcolor{white}{\textbf{Tham số}} & \textcolor{white}{\textbf{Ý nghĩa}} & \textcolor{white}{\textbf{Giá trị khởi đầu}} & \textcolor{white}{\textbf{Tunable?}} \\
\midrule
$\lambda_1$ & Weight Task 1 (Informativeness) & 0.4 & Grid search \\
$\lambda_2$ & Weight Task 2 (Humanitarian) & 0.3 & Grid search \\
$\lambda_3$ & Weight Task 3 (Damage) & 0.3 & Grid search \\
$\alpha$ & Contrastive regularization & 0.1 & Fixed \\
$\beta$ & Graph consistency & 0.05 & Fixed \\
$\gamma$ & Focal Loss focusing param & 2.0 & Ablation \\
\bottomrule
\end{tabular}
\end{table}

\subsection{Adaptive $\lambda$ cho DiffAttn}

Thay $\lambda$ cố định bằng input-dependent:
\begin{equation}
\lambda(x) = \sigma\!\left(\text{MLP}\!\left(\text{concat}(f_I^{\text{graph}}, f_T^{\text{graph}})\right)\right)
\end{equation}

Điểm mới: $\lambda$ nhận \textbf{graph-enhanced features} thay vì raw CLIP features $\rightarrow$ context-aware noise cancellation.

% ================================================================
% SECTION 3: KẾ HOẠCH THỰC NGHIỆM
% ================================================================
\section{Kế Hoạch Thực Nghiệm}

\subsection{Experiment 1: Reproduce Baselines + GEDA Comparison (MT1, MT5)}

\begin{taskbox}[Exp 1 — Baseline Reproduction \& GEDA Comparison]
\textbf{Mục tiêu:} Reproduce cả 2 baselines, sau đó so sánh với GEDA trên cùng framework\\
\textbf{Trạng thái:} \textcolor{done}{$\checkmark$ Evaluation framework} $\cdot$ \textcolor{done}{$\checkmark$ Colab notebook} $\cdot$ \textcolor{inprog}{$\triangleright$ Chạy baselines} $\cdot$ \textcolor{todo}{$\square$ GEDA}\\[0.3cm]
\begin{tabular}{llll}
\toprule
\textbf{Model} & \textbf{Protocol} & \textbf{Report} & \textbf{Status} \\
\midrule
Paper 1 (GNN) & 4 sizes $\times$ 10 splits = 40 runs & wF1, bAcc $\pm$ std, CI & \textcolor{inprog}{Sẵn sàng chạy} \\
Paper 2 (DiffAttn) & 3 tasks $\times$ 5 seeds = 15 runs & wF1, Acc $\pm$ std, CI & \textcolor{inprog}{Sẵn sàng chạy} \\
\textbf{GEDA (ours)} & 3 tasks $\times$ 5 seeds $\times$ 4 sizes & All metrics $\pm$ std, CI & \textcolor{todo}{Chờ implement} \\
\bottomrule
\end{tabular}
\\[0.3cm]
\textbf{Đã hoàn thành:}
\begin{itemize}[nosep]
  \item \textcolor{done}{$\checkmark$} Evaluation framework: CI, paired t-test, Bonferroni, Cohen's d
  \item \textcolor{done}{$\checkmark$} Colab notebook \texttt{GEDA\_Baseline\_Comparison.ipynb}: 8 phases tự động
  \item \textcolor{done}{$\checkmark$} Gap analysis: sửa 9 lỗi tham số Paper 1 (arch, fusion, epochs...)
  \item \textcolor{done}{$\checkmark$} Runner script \texttt{geda\_colab\_runner.py}: pipeline tổng quát
\end{itemize}
\textbf{Output:} \texttt{all\_results.csv} + plots + LaTeX tables tự động
\end{taskbox}

\subsection{Experiment 2: GEDA Architecture Ablation (MT1)}

\begin{taskbox}[Exp 2 — Ablation Study]
\textbf{Mục tiêu:} Chứng minh mỗi component đóng góp giá trị\\[0.3cm]
\begin{tabular}{clccc}
\toprule
\rowcolor{tableheader}
\textcolor{white}{\#} & \textcolor{white}{\textbf{Variant}} & \textcolor{white}{\textbf{Graph}} & \textcolor{white}{\textbf{Attention}} & \textcolor{white}{\textbf{MTL}} \\
\midrule
A1 & CLIP + Linear (baseline) & $\times$ & $\times$ & $\times$ \\
A2 & CLIP + GNN only (Paper 1) & \checkmark & $\times$ & $\times$ \\
A3 & CLIP + Attention only (Paper 2) & $\times$ & \checkmark & $\times$ \\
A4 & GNN + Self-Attn & \checkmark & SA & $\times$ \\
A5 & GNN + SA + GCA & \checkmark & SA+GCA & $\times$ \\
A6 & GNN + SA + GCA + DiffAttn & \checkmark & Full & $\times$ \\
A7 & \textbf{GEDA Full (A6 + MTL)} & \checkmark & Full & \checkmark \\
\bottomrule
\end{tabular}
\\[0.3cm]
\textbf{Expected:} A1 $<$ A2 $\approx$ A3 $<$ A4 $<$ A5 $<$ A6 $<$ A7
\end{taskbox}

\subsection{Experiment 3: FAISS Scalability (MT2)}

\begin{taskbox}[Exp 3 — FAISS vs Brute-force]
\textbf{Mục tiêu:} Chứng minh FAISS giữ accuracy trong khi giảm complexity\\[0.3cm]
\begin{tabular}{llll}
\toprule
\textbf{Method} & \textbf{Recall@16} & \textbf{Memory} & \textbf{Build Time} \\
\midrule
Brute-force cosine & 100\% (exact) & $O(n^2)$ & baseline \\
FAISS IVF (nlist=100) & Target: $\geq$95\% & $O(n)$ & measure \\
FAISS HNSW & Target: $\geq$98\% & $O(n \log n)$ & measure \\
\bottomrule
\end{tabular}
\\[0.3cm]
\textbf{Acceptance:} Recall@16 $\geq$ 95\% AND F1 drop $<$ 1\%
\end{taskbox}

\subsection{Experiment 4: Few-shot Performance (MT1)}

\begin{taskbox}[Exp 4 — Few-shot Comparison]
So sánh GEDA vs baselines trên 4 settings:
\\[0.3cm]
\begin{tabular}{lccccc}
\toprule
\rowcolor{tableheader}
\textcolor{white}{\textbf{Model}} & \textcolor{white}{\textbf{50}} & \textcolor{white}{\textbf{100}} & \textcolor{white}{\textbf{250}} & \textcolor{white}{\textbf{500}} & \textcolor{white}{\textbf{Avg}} \\
\midrule
Paper 1 (GNN) & $\square$ & $\square$ & $\square$ & $\square$ & $\square$ \\
Paper 2 (DiffAttn) & $\square$ & $\square$ & $\square$ & $\square$ & $\square$ \\
\textbf{GEDA (ours)} & $\square$ & $\square$ & $\square$ & $\square$ & $\square$ \\
\bottomrule
\end{tabular}
\\[0.3cm]
$\times$ 3 tasks $\times$ 5 seeds = \textbf{180 experiment runs} tổng cộng\\
\textbf{Statistical test:} Paired t-test (p $<$ 0.05) với Bonferroni correction
\end{taskbox}

\subsection{Experiment 5: Cross-Disaster-Type (MT4)}

\begin{taskbox}[Exp 5 — Leave-one-TYPE-out]
\textbf{Mục tiêu:} Đánh giá generalization across natural disaster types
\\[0.3cm]
\begin{tabular}{llll}
\toprule
\rowcolor{tableheader}
\textcolor{white}{\textbf{Exp}} & \textcolor{white}{\textbf{Train Types}} & \textcolor{white}{\textbf{Test Type}} & \textcolor{white}{\textbf{F1}} \\
\midrule
5A & Earthquake + Wildfire + Flood & \textbf{Hurricane} & $\square$ \\
5B & Hurricane + Wildfire + Flood & \textbf{Earthquake} & $\square$ \\
5C & Hurricane + Earthquake + Flood & \textbf{Wildfire} & $\square$ \\
5D & Hurricane + Earthquake + Wildfire & \textbf{Flood} & $\square$ \\
\bottomrule
\end{tabular}
\\[0.3cm]
\textbf{Bonus:} So sánh GEDA vs Paper 1 vs Paper 2 trên cùng splits
\end{taskbox}

% ================================================================
% SECTION 4: DELIVERABLES
% ================================================================
\section{Sản Phẩm Đầu Ra (Deliverables)}

\begin{table}[H]
\centering
\caption{Danh sách deliverables theo phase}
\rowcolors{2}{tablerow1}{tablerow2}
\begin{tabularx}{\textwidth}{c l X c}
\toprule
\rowcolor{tableheader}
\textcolor{white}{\#} & \textcolor{white}{\textbf{Deliverable}} & \textcolor{white}{\textbf{Mô tả}} & \textcolor{white}{\textbf{Phase}} \\
\midrule
D1 & \texttt{baseline\_5seeds.csv} & Kết quả 2 baselines, 5 seeds, 3 tasks, 4 settings & P1 \\
D2 & \texttt{geda\_model.py} & Code kiến trúc GEDA (PyTorch) & P2 \\
D3 & \texttt{faiss\_graph.py} & Module graph construction với FAISS & P2 \\
D4 & \texttt{mtl\_trainer.py} & Multi-task training loop + Focal Loss & P2 \\
D5 & \texttt{ablation\_results.csv} & Kết quả 7 variants ablation study & P2 \\
D6 & \texttt{cross\_type\_results.csv} & Kết quả 4 leave-one-type-out experiments & P3 \\
D7 & \texttt{geda\_notebook.ipynb} & Colab notebook chạy end-to-end & P3 \\
D8 & Bài báo nghiên cứu & LaTeX paper submission-ready & P4 \\
D9 & Code release & GitHub repo public + README & P4 \\
\bottomrule
\end{tabularx}
\end{table}

% ================================================================
% SECTION 5: TIMELINE
% ================================================================
\section{Timeline}

\subsection{Tổng Quan 4 Phases}

\begin{figure}[H]
\centering
\begin{tikzpicture}[
  phase/.style={rectangle, rounded corners=4pt, minimum width=3.5cm, minimum height=1.5cm, align=center, font=\small\bfseries, text=white},
  arrow/.style={-{Stealth[length=3mm]}, thick, color=gray!60},
  label/.style={font=\tiny, color=gray}
]
  \node[phase, fill=phase1] (p1) {Phase 1\\Foundations\\2--3 tuần};
  \node[phase, fill=phase2, right=0.8cm of p1] (p2) {Phase 2\\Core GEDA\\3--4 tuần};
  \node[phase, fill=phase3, right=0.8cm of p2] (p3) {Phase 3\\Evaluation\\2--3 tuần};
  \node[phase, fill=phase4, right=0.8cm of p3] (p4) {Phase 4\\Paper\\2--3 tuần};
  
  \draw[arrow] (p1) -- (p2);
  \draw[arrow] (p2) -- (p3);
  \draw[arrow] (p3) -- (p4);
  
  \node[label, below=0.2cm of p1] {Tuần 1--3};
  \node[label, below=0.2cm of p2] {Tuần 4--7};
  \node[label, below=0.2cm of p3] {Tuần 8--10};
  \node[label, below=0.2cm of p4] {Tuần 11--13};
\end{tikzpicture}
\caption{4 phases của dự án ($\sim$13 tuần)}
\end{figure}

\subsection{Chi Tiết Từng Phase}

\subsubsection{Phase 1: Foundations (Tuần 1--3)}

\begin{table}[H]
\centering
\rowcolors{2}{tablerow1}{tablerow2}
\begin{tabularx}{\textwidth}{c X c c c}
\toprule
\rowcolor{phase1!90}
\textcolor{white}{\textbf{\#}} & \textcolor{white}{\textbf{Task}} & \textcolor{white}{\textbf{Tuần}} & \textcolor{white}{\textbf{Output}} & \textcolor{white}{\textbf{Status}} \\
\midrule
1.1 & Reproduce Paper 1 (GNN) -- 4 sizes $\times$ 10 splits & 1 & D1 & \textcolor{inprog}{$\triangleright$} \\
1.2 & Reproduce Paper 2 (DiffAttn) -- 3 tasks $\times$ 5 seeds & 1 & D1 & \textcolor{inprog}{$\triangleright$} \\
1.3 & Setup evaluation framework (CI, paired t-test, Cohen's d) & 1 & Script & \textcolor{done}{$\checkmark$} \\
1.4 & Implement Focal Loss + class-weighted sampling & 2 & Module & \textcolor{todo}{$\square$} \\
1.5 & FAISS benchmark: Recall@k vs brute-force & 2 & D3 & \textcolor{todo}{$\square$} \\
1.6 & Chuẩn bị CrisisMMD data loaders cho MTL & 2--3 & Module & \textcolor{todo}{$\square$} \\
1.7 & Viết unit tests cho graph construction & 3 & Tests & \textcolor{todo}{$\square$} \\
\bottomrule
\end{tabularx}
\end{table}

\subsubsection{Phase 2: Core GEDA Implementation (Tuần 4--7)}

\begin{table}[H]
\centering
\rowcolors{2}{tablerow1}{tablerow2}
\begin{tabularx}{\textwidth}{c X c c}
\toprule
\rowcolor{phase2!90}
\textcolor{white}{\textbf{\#}} & \textcolor{white}{\textbf{Task}} & \textcolor{white}{\textbf{Tuần}} & \textcolor{white}{\textbf{Output}} \\
\midrule
2.1 & Implement GEDA model: GNN + Attention stack & 4 & D2 \\
2.2 & Implement Adaptive $\lambda(x)$ cho DiffAttn & 4 & D2 \\
2.3 & Implement Multi-task training loop & 5 & D4 \\
2.4 & Hyperparameter tuning ($\lambda_1, \lambda_2, \lambda_3, \gamma, k$) & 5 & Config \\
2.5 & Ablation study: 7 variants $\times$ 5 seeds & 6 & D5 \\
2.6 & Few-shot comparison: GEDA vs baselines & 6--7 & Results \\
2.7 & Debug + optimize training pipeline & 7 & D7 \\
\bottomrule
\end{tabularx}
\end{table}

\subsubsection{Phase 3: Advanced Evaluation (Tuần 8--10)}

\begin{table}[H]
\centering
\rowcolors{2}{tablerow1}{tablerow2}
\begin{tabularx}{\textwidth}{c X c c}
\toprule
\rowcolor{phase3!90}
\textcolor{white}{\textbf{\#}} & \textcolor{white}{\textbf{Task}} & \textcolor{white}{\textbf{Tuần}} & \textcolor{white}{\textbf{Output}} \\
\midrule
3.1 & Cross-disaster-type experiments (4 splits) & 8 & D6 \\
3.2 & Adaptive $\lambda$ analysis: task-specific vs global & 8--9 & Analysis \\
3.3 & Confusion matrix + per-class P/R/F1 & 9 & Tables \\
3.4 & Latency benchmark (training + inference) & 9 & Benchmark \\
3.5 & Error analysis: failure cases visualization & 10 & Figures \\
3.6 & Statistical significance tests cho tất cả comparisons & 10 & Report \\
\bottomrule
\end{tabularx}
\end{table}

\subsubsection{Phase 4: Paper Writing + Release (Tuần 11--13)}

\begin{table}[H]
\centering
\rowcolors{2}{tablerow1}{tablerow2}
\begin{tabularx}{\textwidth}{c X c c}
\toprule
\rowcolor{phase4!90}
\textcolor{white}{\textbf{\#}} & \textcolor{white}{\textbf{Task}} & \textcolor{white}{\textbf{Tuần}} & \textcolor{white}{\textbf{Output}} \\
\midrule
4.1 & Viết paper: Introduction, Related Work, Method & 11 & D8 \\
4.2 & Viết paper: Experiments, Results, Discussion & 12 & D8 \\
4.3 & Tạo figures, tables, diagrams cho paper & 12 & Figures \\
4.4 & Code cleanup + documentation + README & 12--13 & D9 \\
4.5 & Internal review + revision & 13 & D8 final \\
\bottomrule
\end{tabularx}
\end{table}

% ================================================================
% SECTION 6: COMPUTE RESOURCES
% ================================================================
\section{Tài Nguyên Compute}

\begin{table}[H]
\centering
\caption{Ước tính tài nguyên compute}
\rowcolors{2}{tablerow1}{tablerow2}
\begin{tabularx}{\textwidth}{l X c c}
\toprule
\rowcolor{tableheader}
\textcolor{white}{\textbf{Experiment}} & \textcolor{white}{\textbf{Chi tiết}} & \textcolor{white}{\textbf{GPU Hours}} & \textcolor{white}{\textbf{Platform}} \\
\midrule
Baseline reproduce & 2 models $\times$ 5 seeds $\times$ 4 settings $\times$ 3 tasks & $\sim$40h & Colab \\
GEDA ablation & 7 variants $\times$ 5 seeds $\times$ 1 setting & $\sim$35h & Colab \\
Few-shot full & 1 model $\times$ 5 seeds $\times$ 4 settings $\times$ 3 tasks & $\sim$60h & Colab \\
Cross-type & 4 splits $\times$ 5 seeds $\times$ 3 models & $\sim$30h & Colab \\
Hyperparameter & Grid search $\sim$50 configs & $\sim$25h & Colab \\
\midrule
\textbf{Tổng} & & \textbf{$\sim$190h} & \\
\bottomrule
\end{tabularx}
\end{table}

\begin{riskbox}[Lưu ý Compute]
Google Colab Free: $\sim$4h/session. Colab Pro: $\sim$24h/session.\\
Ước tính: $\sim$190 GPU hours $\approx$ \textbf{8 ngày liên tục} (L4 GPU) hoặc $\sim$48 Colab Free sessions.\\
\textbf{Khuyến nghị:} Sử dụng Colab Pro ($\sim$\$10/tháng) để đảm bảo tiến độ.
\end{riskbox}

% ================================================================
% SECTION 7: RISK MANAGEMENT
% ================================================================
\section{Quản Lý Rủi Ro}

\begin{table}[H]
\centering
\caption{Ma trận rủiро và phương án giảm thiểu}
\rowcolors{2}{tablerow1}{tablerow2}
\begin{tabularx}{\textwidth}{c X c c X}
\toprule
\rowcolor{tableheader}
\textcolor{white}{\#} & \textcolor{white}{\textbf{Rủi ro}} & \textcolor{white}{\textbf{P}} & \textcolor{white}{\textbf{I}} & \textcolor{white}{\textbf{Giảm thiểu}} \\
\midrule
R1 & GEDA không outperform baselines & 30\% & High & Ablation sẽ chỉ ra component nào yếu; điều chỉnh \\
R2 & Colab GPU timeout giữa chừng & 50\% & Med & Checkpoint mỗi epoch; resume training \\
R3 & FAISS accuracy drop quá lớn & 20\% & Med & Benchmark trước; fallback brute-force nếu cần \\
R4 & Class imbalance làm metrics lệch & 40\% & High & Focal Loss + class weights + report per-class \\
R5 & CrisisMMD quá nhỏ cho MTL & 35\% & High & Data augmentation + pseudo-labels \\
R6 & Thời gian vượt 13 tuần & 40\% & Med & Ưu tiên MT1 + MT5; cắt bỏ Phase 3.5, 3.6 nếu cần \\
\bottomrule
\end{tabularx}
\end{table}

% ================================================================
% SECTION 8: SUCCESS CRITERIA
% ================================================================
\section{Tiêu Chí Thành Công}

\begin{table}[H]
\centering
\caption{Tiêu chí thành công cho từng mục tiêu}
\rowcolors{2}{tablerow1}{tablerow2}
\begin{tabularx}{\textwidth}{c X X c}
\toprule
\rowcolor{tableheader}
\textcolor{white}{\textbf{MT}} & \textcolor{white}{\textbf{Mục tiêu}} & \textcolor{white}{\textbf{Tiêu chí}} & \textcolor{white}{\textbf{Must/Nice}} \\
\midrule
MT1 & GEDA outperform baselines & F1 cao hơn trên $\geq$ 2/3 tasks, $p < 0.05$ & \textbf{Must} \\
MT2 & Scalability & Recall@16 $\geq$ 95\%, F1 drop $<$ 1\% & Must \\
MT3 & Multi-task Learning & MTL $\geq$ single-task trên $\geq$ 2/3 tasks & Must \\
MT4 & Cross-type generalization & F1 $>$ random baseline trên held-out type & Nice \\
MT5 & Statistical rigor & 5 seeds, CI, sig. tests cho mọi results & \textbf{Must} \\
\bottomrule
\end{tabularx}
\end{table}

\begin{goalbox}[Minimum Viable Research (MVR)]
Nếu thời gian hạn chế, \textbf{MVR} bao gồm:
\begin{enumerate}
  \item Baselines reproduce với 5 seeds (Exp 1) — \textbf{Bắt buộc}
  \item GEDA implementation + ablation (Exp 2) — \textbf{Bắt buộc}
  \item Few-shot comparison trên CrisisMMD (Exp 4) — \textbf{Bắt buộc}
  \item Statistical analysis (CI + tests) — \textbf{Bắt buộc}
\end{enumerate}
Exp 3 (FAISS) và Exp 5 (Cross-type) có thể defer nếu cần.
\end{goalbox}

% ================================================================
% SECTION 9: NEXT STEPS
% ================================================================
\section{Bước Tiếp Theo Ngay Lập Tức}

\begin{enumerate}[leftmargin=*]
  \item \textcolor{done}{$\checkmark$} \textbf{DONE:} Setup evaluation framework
    \begin{itemize}
      \item Script tính CI, paired t-test, Bonferroni, Cohen's d
      \item Template CSV thống nhất cho kết quả
      \item Colab notebook + runner script đã sẵn sàng
      \item Gap analysis: sửa 9 lỗi tham số Paper 1
    \end{itemize}
  \item \textcolor{inprog}{$\triangleright$} \textbf{ĐANG LÀM:} Chạy baselines trên Colab
    \begin{itemize}
      \item Upload \texttt{GEDA\_Baseline\_Comparison.ipynb} lên Colab
      \item Chạy Phase 1--5: Paper 2 (15 runs) + Paper 1 (40 runs)
      \item Phase 6: Phân tích kết quả, so sánh với paper reported values
    \end{itemize}
  \item \textcolor{todo}{$\square$} \textbf{TIẾP THEO:} Implement GEDA model
    \begin{itemize}
      \item Xây dựng \texttt{geda\_model.py}: GNN $\rightarrow$ Self-Attn $\rightarrow$ GCA $\rightarrow$ DiffAttn
      \item Thêm GEDA vào notebook Phase 7 để so sánh trực tiếp
      \item Chạy 3 tasks $\times$ 5 seeds $\times$ 4 few-shot sizes
    \end{itemize}
  \item \textcolor{todo}{$\square$} \textbf{SAU ĐÓ:} Ablation study (7 variants $\times$ 5 seeds)
\end{enumerate}

\vspace{0.5cm}

\begin{goalbox}[Cam kết]
Dự án nghiên cứu này hướng tới đóng góp \textbf{novelty chính}: kết hợp noise handling ở \textit{instance-level} (graph propagation) và \textit{feature-level} (differential attention) --- hai cơ chế bổ trợ nhau mà chưa paper nào khai thác, tập trung hoàn toàn vào bài toán \textbf{phân loại thảm hoạ thiên nhiên đa phương thức} có tính ứng dụng cao trong emergency management.
\end{goalbox}

\end{document}
